\chapter{Limitations and Future Work}
\label{sec:limitations}

\section{Limitations}

\textbf{Untrimmed NGT recordings.} The NGT training videos were not
trimmed: each begins with a lengthy idle segment in which the signer
stands watching the reference video before starting to sign, and ends
with a reach toward the foot pedal used to stop the recording. Both add
frames unrelated to signing --- idle standing and extraneous motion ---
and because the idle portion differs in length per signer, no trivial
automated trim existed; we judged manual cleaning not worth thesis time.
The reported NGT numbers therefore include this noise floor.
\textbf{Frame-independent diffusion editing introduces flicker.} Our
Flux9B pipeline (\S\ref{sec:setup-generation}\todome{ref re-pointed
03.08: pipeline instance detail moved to Setup; was sec:method-diffusion}) edits each frame of a
source video independently: every edited frame follows the appearance-edit
prompt on its own, with no explicit mechanism enforcing consistency with
the frame before or after it, so consecutive edited frames drift slightly
relative to one another and the resulting clip flickers. We looked for a
video-native alternative that would edit a clip as one temporally
consistent unit rather than frame-by-frame, trying LTX 2.3 (22B)
~\cite{hacohen2026ltx2}, LTX 2.3 with an Edit Anything
LoRA~\cite{alissonerdx2026editanything}, Wan
2.2~\cite{wan2025}, and Lucy Edit 1.1 dev
(5B)~\cite{decart2025lucyedit}\todome{citations added 04.08,
web-verified; caveats: the Wan cite is the team's own suggested citation
but is technically the 2.1 tech report (no 2.2 paper exists); LTX-2.3
has no own paper, Lightricks points to the LTX-2 report;
Edit-Anything is a community LoRA (no paper) cited as @misc to its HF
repo --- drop that cite and keep just the name if you prefer};
in every case the edited
output was substantially worse than the frame-by-frame Flux9B route, so we
kept the per-frame pipeline and report the resulting flicker as a
limitation rather than a solved problem.

\textbf{Occasional diffusion hallucinations.} The diffusion model
occasionally hallucinates content that was not in the source frame ---
most visibly, adding a spurious or malformed handshape to a frame whose
original signer was not making that shape, or to a frame that did not
contain a visible hand at all. This risk is highest for the skin-tone and
signer-swap edit types (\S\ref{sec:method-diffusion}), and is a direct
consequence of editing each frame independently rather than conditioning
on the sign actually being performed.
\todo{example frames:  source frames from nano\_banana/output/flux-aug} A more heavily
constrained editing pipeline --- for instance, conditioning the editor
on the source pose via ControlNet --- could reduce these hallucinations
further, but their prevalence was not high enough to justify the added
complexity.\todome{sentence added 04.08 from your todo}

\section{Future work}
\label{sec:future-work}
\todo{future work is too long, decide what is actually important. Make things less verbose}
\textbf{Parametric avatars in place of per-avatar mocap retargeting.} The
natural next step for the NGT pipeline is to replace hand-built
CharacterCreator avatars with parametric SMPL-X
bodies~\cite{pavlakosExpressiveBodyCapture2019a}: the same Unreal rendering
strengths, but with appearance and body shape as sampled parameters rather
than authored assets, enabling appearance sweeps at scale. Diffusion-based
sign production directly in SMPL-X pose
space~\cite{baltatzisNeuralSignActors2024a} and large 3D sign
corpora~\cite{yuSignAvatarsLargeScale3D2025} make this increasingly
practical --- though driving avatars from estimated or generated SMPL-X
parameters remains harder than working directly with recorded motion
capture, as our pipeline does.

\textbf{Signer identity persists through the synthetic-avatar swap.}
\todome{moved here from Limitations 04.08 per your todo; the second half
of that todo (``make sure it is also addressed in the discussion'') is
still open --- the natural spot is the identity-diversity attribution in
\S\ref{sec:results-crossdomain}/the regime discussion, which currently
credits the NGT gain to identity diversity without acknowledging that
the avatars retain the source signer's motion signature} Even
when signer identity is changed by re-rendering with a synthetic avatar,
the avatar does not fully erase who provided the underlying motion
capture: body shape, body posture, range of motion when signing, and
signing speed are all retargeted along with the sign itself, and all four
remain visibly present in the render. This is not only a probe result
(\S\ref{sec:results-diagnosis}) but a plain visual one: as someone with no
knowledge of sign language, when looking at individual frames of our
synthetic-signer renders I can still tell which of the real NGT signers
supplied the motion capture for a given avatar. The synthetic-signer
intervention changes surface appearance --- body model, clothing, skin
tone --- without fully removing the motion-level signature of the source
signer, a distinction the retrieval-based evaluation in this thesis does
not separate out. Measuring that motion-level signature directly --- a
signer-identity probe on the rendered views themselves --- would be a
cheap first step, and motion retargeting that also perturbs kinematic
style the harder follow-up.

\textbf{Scaling laws for synthetic signers.} The synthetic-signer
dose-response (0/1/2 avatars, and beyond with new renders) remains the
most valuable follow-up on the synthetic axis: does the NGT gain grow,
saturate, or invert as synthetic appearance diversity increases?

\textbf{Augmenting the full ASL Citizen training set, not just a targeted
subset.} Our diffusion augmentation covers a deliberately targeted subset
of the ASL Citizen training set, chosen worst-first by signer scarcity and
baseline difficulty (\S\ref{sec:method-diffusion}) --- a choice guided in
part by evidence that targeted synthetic augmentation at smaller scale can
outperform augmenting everything~\cite{nguyenDoWeNeedSynthetic2025}.
Whether that finding transfers to our setting is untested: it remains an
open question whether augmenting the entire training set, rather than the
selected worst-first subset, would have produced a larger in-domain gain
than the targeted budget did, or whether the targeted-subset advantage
holds here too and full-corpus augmentation would mostly add compute cost
without adding retrieval gain. The compute constraint that motivated the
original subset choice (\S\ref{sec:setup-generation}\todome{ref
re-pointed 03.08: the compute-budget rationale moved to Setup; was
sec:method-diffusion}) would need to be
relaxed to test this directly.

\textbf{End-to-end training for the NGT encoder.} All NGT results come
from an encoder trained on pre-extracted, frozen Logos clip features
(\S\ref{sec:results-crossdomain}). An end-to-end variant that prepends
the MViTv2-S backbone itself --- a frozen-backbone control and a fully
fine-tuned configuration --- was implemented and configured but never
run. Running both would show whether the cross-domain gain from
synthetic views survives, grows, or shrinks once gradients can reach the
visual backbone, giving the low-data regime its own counterpart to the
in-domain contrast between frozen-feature heads and backbone
fine-tuning (\S\ref{sec:discussion-regimes}).

\textbf{A discriminative term for the high-$K$ arms.} \todo{can we just run it now since we have more compute?} In the NGT
ladder, retrieval falls as more counterfactual views are packed into
each training item ($K{\geq}5$ arms below their $K{=}3$ counterparts),
and the pair-SupCon head points the same way on ASL Citizen ($0.567$ at
$K{=}3$, against $0.800$ for the same data without explicit pairing);
the convergence audit rules out undertraining.\todome{rephrased 04.08 to
drop the removed K=5 retrieval number; original cited K=3-to-K=5
0.567-to-0.414} Our reading is that pure
SupCon at high $K$ spends its capacity making views of the same item
indistinguishable, at the cost of keeping different items apart. A
natural test is to add a discriminative term to the NGT training setup
(a 12-layer encoder trained with SupCon on frozen Logos clip features;
\S\ref{sec:results-crossdomain}): cross-entropy on sentence
(or gloss) identity plus $\lambda$ times the SupCon term, sweeping
$\lambda \in \{0.1, 1\}$ over a subset of arms (E1, E3, E5, E8, 3 seeds
each, ${\approx}700$ GPU-hours-equivalent SBU). If the high-$K$ drop
disappears, the tradeoff is objective-specific and fixable; if it
persists, appearance invariance itself is what costs retrieval.\todo{framing check}

\textbf{Appearance-matched hard negatives.} Our pair objectives use the
synthetic edits as appearance-matched \emph{positives}: views that
change the signer while preserving the motion are pulled toward their
source clip. The mirror image is to use the same augmentation for hard
\emph{negatives}: clips sharing one synthetic identity (the same avatar,
or the same Flux signer edit) but performing different sentences or
glosses, so that appearance can no longer separate the negative pair and
only motion can. This drops into the existing SupCon setup as a
batch-construction change rather than a new loss: negatives are
currently simply the other items in the batch
(\S\ref{sec:method-objectives}), and would instead be sampled so that
each anchor also meets its own synthetic identity performing other
signs. The expected insight bears on the form-dose result: explicit
motion-discrimination pressure would test whether the form organisation
that pair-trained heads currently \emph{amplify}
(\S\ref{sec:results-iconicity}, Fig.~\ref{fig:form_dose}) recedes once
negatives, and not only positives, are appearance-controlled --- or
sharpens further as the space reorganises around
motion.\todo{framing check}

\textbf{A from-scratch control for the in-domain null.} \todo{refine writing, this is the experiment suggested by Angelos} The two-regime
reading attributes the flat in-domain result to invariance already
consumed during pretraining (\S\ref{sec:discussion-regimes}), and this
is directly testable: train the same MViTv2-S from random initialisation
on ASL Citizen, with cross-entropy alone and with the synthetic
augmentations added, and compare against the pretrained baseline. If the
from-scratch model lands well below the pretrained one and the
augmentations now help it, pretraining diversity --- not a saturated
recipe --- explains the null; if augmentation still adds nothing, the
in-domain ceiling is not a pretraining artifact. 

\textbf{Fine-tuning I3D as a second backbone for the embedding-space
analyses.} Throughout this thesis I3D is used only as a frozen,
off-the-shelf feature extractor (\S\ref{sec:results-diagnosis},
\S\ref{sec:results-indomain}), never fine-tuned the way the Logos/MViTv2-S
backbone is. Repeating the fine-tuning runs and the embedding-space
analyses --- signer-separability probes, form-dose geometry --- with a
fine-tuned I3D would show whether this thesis's diagnostic findings are
specific to the MViTv2-S/Logos backbone or hold across RGB architectures
more generally.
